\chapter{Research Plan}
\label{chap:plan}

\section{Expected contribution}

The dissertation will defend the hypothesis of Chapter~\ref{chap:intro} and deliver a reusable artifact (Figure~\ref{fig:pipeline}):
\begin{enumerate}
  \item a small one-shot selector, fine-tuned for multi-hop evidence, with adaptive stopping;
  \item a \emph{chain adapter}: a rank-2 update (about 180,000 parameters, under 1\% of the model) that turns the one-shot selector into an iterative one. A single base model then serves both modes;
  \item a router that decides, per question, whether to activate the adapter;
  \item an evaluation script that reports tokens, passes, full-support recall and answer quality against fixed top-$k$, Adaptive-$k$ and always-iterative baselines.
\end{enumerate}
The adapter saves memory and deployment complexity, not passes: activating it still costs iterative inference. Savings in passes depend on the router.

\begin{figure}[htbp]
\centering
\resizebox{\textwidth}{!}{% Planned system: cheap selector, router and rank-2 chain adapter (TikZ).
\begin{tikzpicture}[
  box/.style={draw, rounded corners, align=center, minimum height=1.0cm, font=\small},
  arr/.style={-{Latex[length=2mm]}, thick},
  node distance=0.6cm and 0.9cm]
  \node[box, minimum width=2.2cm] (q) {question +\\$\sim$20 paragraphs};
  \node[box, right=of q, fill=gray!10, minimum width=2.8cm] (sel) {small selector\\one-shot, 20 passes};
  \node[box, right=of sel, minimum width=2.0cm] (rt) {router};
  \node[box, above right=0.4cm and 0.9cm of rt, fill=orange!15, minimum width=3.2cm] (ad) {selector + chain adapter\\(rank 2), iterative};
  \node[box, below right=0.4cm and 0.9cm of rt, minimum width=3.2cm] (keep) {keep one-shot ranking};
  \node[box, right=4.6cm of rt, fill=blue!8, minimum width=2.4cm] (llm) {adaptive stop\\$\rightarrow$ LLM reader};
  \draw[arr] (q) -- (sel);
  \draw[arr] (sel) -- (rt);
  \draw[arr] (rt) |- node[pos=0.75, above, font=\footnotesize] {long / hard} (ad);
  \draw[arr] (rt) |- node[pos=0.75, below, font=\footnotesize] {otherwise} (keep);
  \draw[arr] (ad) -| (llm);
  \draw[arr] (keep) -| (llm);
\end{tikzpicture}
}
\caption{Planned system. The cheap one-shot selector runs for every question; the router activates the rank-2 chain adapter only where iteration is expected to help; adaptive stopping decides how many paragraphs reach the LLM.}
\label{fig:pipeline}
\end{figure}

\section{Remaining work}

\begin{description}
  \item[W1. Train the adapter directly.] So far the rank-2 change was extracted from a full fine-tuning run. Training a rank-1, 2 and 4 LoRA \cite{hu2022lora} on the same pairs will show whether the artifact can be built directly, and how it compares with the extracted one \cite{biderman2024loralearnsless,shuttleworth2025illusion}.
  \item[W2. Router trained on benefit.] The hop-count router failed (Section~\ref{sec:res_long}). A router trained to predict, per question, whether iteration improves the top five is pre-registered and waiting for GPU time.
  \item[W3. Answer quality.] Whether the recall gain on long chains reaches the reader's answers (Section~\ref{sec:res_f1}).
  \item[W4. Geometry with noise baselines.] Singular-value energy, stable rank and principal angles of the conditioning change, compared with seed-to-seed differences and random matrices \cite{altaf2026noisebaseline}, to separate functional rank from spectral structure.
  \item[W5. Transfer.] Graft the adapter onto other one-shot rerankers (A1, the off-the-shelf MS MARCO model) and evaluate on 2WikiMultiHopQA \cite{ho2020twowiki}. Transfer of a behaviour, not only a domain, between rerankers would be the main novelty \cite{braga2025taskarithmetic}.
\end{description}

\section{Threats to validity}

\begin{itemize}
  \item \textbf{Effect size.} The confirmed gain on long chains is modest ($+0.055$ recall@5 on four-hop questions) and heterogeneous between data sources.
  \item \textbf{Single architecture and data source.} All selectors are MiniLM-L6 and all evaluation questions come from MuSiQue (and HotpotQA for the first results). W5 addresses this.
  \item \textbf{Evaluation sets.} MuSiQue test labels are not public; evaluation uses disjoint subsets of the development split, and one fresh set has been used twice (with the second use declared in advance).
  \item \textbf{Multiplicity.} Each experiment has a single primary test; across experiments no correction is applied. The headline claims will be re-tested on a locked set before the defence.
\end{itemize}

\section{Schedule}

Table~\ref{tab:schedule} lists the planned activities by quarter.

\begin{table}[htbp]
\centering
\caption{Schedule (quarters from the qualification).}
\label{tab:schedule}
\small
\begin{tabular}{lcccc}
\toprule
Activity & Q1 & Q2 & Q3 & Q4 \\
\midrule
W1 adapter trained directly; W3 answer quality & $\bullet$ & & & \\
W2 benefit router and cascade & $\bullet$ & $\bullet$ & & \\
W4 geometry with noise baselines & & $\bullet$ & & \\
W5 transfer across rerankers and to 2WikiMultiHopQA & & $\bullet$ & $\bullet$ & \\
Paper submission & & & $\bullet$ & \\
Locked-set confirmation of the headline claims & & & $\bullet$ & \\
Dissertation writing and defence & & & $\bullet$ & $\bullet$ \\
\bottomrule
\end{tabular}
\end{table}
